\section{Theory and hypotheses}\label{sec:theory}

Three literatures jointly motivate the analysis: performance theories of
political trust, structural-grievance accounts of dislocation, and the
empirical occupation-exposure tradition. Each contributes a distinct
prediction about the direction of the AI-exposure effect on
institutional trust, and the three predictions partly conflict. The
within--between specification used here is what allows the data to
adjudicate between them.

\paragraph{Institutional performance and trust.}
\citet{Hetherington1998} and \citet{vanderMeer2017} argue that political
trust responds to \emph{perceived institutional performance}, including
how well institutions appear to manage macroeconomic and structural
shocks. A technological shock that visibly weakens an institution's
capacity to deliver economic security---through job destruction, sectoral
churn, or obsolete policy instruments---should therefore lower trust on
average. The account is symmetric: where active labour-market policy,
retraining, or strong dismissal protection appear to absorb the shock,
trust can hold or even rise.

\paragraph{Structural grievance.}
\citet{GidronHall2017} and \citet{Kurer2020} sharpen the performance
account into a compositional mechanism: technological change creates
labour-market losers, and losers re-route their political behaviour
through grievance. The implication for trust is individual-level:
\emph{respondents} in occupations more exposed to substitution should
report lower trust, conditional on country context.

\paragraph{Exposure measures and political behaviour.}
The empirical occupation-exposure literature is dominated by
within-country designs. Trade-shock studies use regional import
competition (\citealt{ColantoneStanig2018} for Brexit). Automation
studies use individual robot exposure
(\citealt{AnelliColantoneStanig2021} for Western Europe;
\citealt{FreyBergerChen2018} for U.S.\ counties). \citet{AcemogluRestrepo2020}
ground the labour-market mechanism in U.S.\ wage and employment data.
These designs identify effects from \emph{between-region within-country}
variation in exposure---the within-country margin of a multilevel model.
What they do not show is whether the cross-national pattern carries the
same sign.

\paragraph{Contribution.}
This paper extends the tradition along three margins. \emph{First}, it
substitutes a manufacturing-biased robot measure for the
white-collar-leaning ILO--NASK GenAI exposure score \citep{Gmyrek2025},
which is ISCO-08 4-digit native and carries a 2023$\to$2025
capability-vintage shift. Because that shift is common to all countries,
it is absorbed by the round dummies and does not by itself identify the
within-country slope; the within variation that survives detrending is
the smaller, country-specific occupational recomposition.
\emph{Second}, it pools cross-national data and applies the within--between
specification \citep{BellJones2015,SchmidtCatranFairbrother2016} so that
the between- and within-country slopes are estimated separately rather
than constrained to be equal. \emph{Third}, it pre-specifies an
education-moderation test with two competing signed predictions and a
precise null, treating a tight null as theoretically informative.

\paragraph{Hypotheses.}
\begin{description}
\setlength\itemsep{2pt}
  \item[H1 (between).] Across countries, higher mean AI exposure
    $\bar E_c$ predicts lower mean institutional trust:
    $\gamma_B < 0$. Predicted by structural grievance if between-country
    variation in exposure tracks the size of the loser pool.
  \item[H2 (within).] Within countries, increases in exposure
    $(E_{ct}-\bar E_c)$ erode trust: $\gamma_W < 0$. Predicted by the
    performance account when institutions visibly fail to absorb the
    shock.
  \item[H2a (decomposition).] $\gamma_W \neq \gamma_B$: the
    cross-sectional and longitudinal margins answer different questions.
    Tested via the Mundlak Wald.
  \item[H3 (compositional).] Individuals in AI-exposed occupations
    report lower trust net of country context:
    $\beta_{\mathrm{genai}} < 0$. The classic structural-grievance
    prediction at the individual level.
  \item[H4.] The within-country exposure effect is moderated by
    education. Direction is theoretically underspecified:
    H4a (substitution) predicts a \emph{negative} interaction
    (higher-educated workers react more, since GenAI substitutes their
    cognitive tasks); H4b (curator) predicts a \emph{positive} one
    (higher-educated workers move into oversight and curator roles
    that complement AI); H4$_0$ predicts no moderation.
  \item[H5 (institutional moderation).] Stronger employment-protection
    legislation attenuates the within-effect:
    \emph{positive} coefficient on $\mathrm{EPL}_c \times (E_{ct}-\bar E_c)$.
\end{description}
